\documentclass[10pt,a4paper]{amsart}
\usepackage{amsmath,amssymb,amsthm,graphicx}
\usepackage{verbatim}
\newenvironment{reference}{\comment}{\endcomment}
%\newenvironment{reference}{}{}
\newenvironment{slogan}{\comment}{\endcomment}
\newenvironment{history}{\comment}{\endcomment}

% For commutative diagrams we use Xy-pic
\usepackage[all]{xy}

% We use 2cell for 2-commutative diagrams.
\xyoption{2cell}
% Theorem environments.
%
\theoremstyle{plain}
\newtheorem{theorem}[subsection]{Theorem}
\newtheorem{proposition}[subsection]{Proposition}
\newtheorem{lemma}[subsection]{Lemma}

\theoremstyle{definition}
\newtheorem{definition}[subsection]{Definition}
\newtheorem{example}[subsection]{Example}
\newtheorem{exercise}[subsection]{Exercise}
\newtheorem{situation}[subsection]{Situation}

\theoremstyle{remark}
\newtheorem{remark}[subsection]{Remark}
\newtheorem{remarks}[subsection]{Remarks}

\numberwithin{equation}{subsection}

% Macros
%
\def\lim{\mathop{\mathrm{lim}}\nolimits}
\def\colim{\mathop{\mathrm{colim}}\nolimits}
\def\Spec{\mathop{\mathrm{Spec}}}
\def\Hom{\mathop{\mathrm{Hom}}\nolimits}
\def\Ext{\mathop{\mathrm{Ext}}\nolimits}
\def\SheafHom{\mathop{\mathcal{H}\!\mathit{om}}\nolimits}
\def\SheafExt{\mathop{\mathcal{E}\!\mathit{xt}}\nolimits}
\def\Sch{\mathit{Sch}}
\def\Mor{\mathop{\mathrm{Mor}}\nolimits}
\def\Ob{\mathop{\mathrm{Ob}}\nolimits}
\def\Sh{\mathop{\mathit{Sh}}\nolimits}
\def\NL{\mathop{N\!L}\nolimits}
\def\CH{\mathop{\mathrm{CH}}\nolimits}
\def\proetale{{pro\text{-}\acute{e}tale}}
\def\etale{{\acute{e}tale}}
\def\QCoh{\mathit{QCoh}}
\def\Ker{\mathop{\mathrm{Ker}}}
\def\Im{\mathop{\mathrm{Im}}}
\def\Coker{\mathop{\mathrm{Coker}}}
\def\Coim{\mathop{\mathrm{Coim}}}

% Boxtimes
%
\DeclareMathSymbol{\boxtimes}{\mathbin}{AMSa}{"02}

%
% Macros for moduli stacks/spaces
%
\def\QCohstack{\mathcal{QC}\!\mathit{oh}}
\def\Cohstack{\mathcal{C}\!\mathit{oh}}
\def\Spacesstack{\mathcal{S}\!\mathit{paces}}
\def\Quotfunctor{\mathrm{Quot}}
\def\Hilbfunctor{\mathrm{Hilb}}
\def\Curvesstack{\mathcal{C}\!\mathit{urves}}
\def\Polarizedstack{\mathcal{P}\!\mathit{olarized}}
\def\Complexesstack{\mathcal{C}\!\mathit{omplexes}}
% \Pic is the operator that assigns to X its picard group, usage \Pic(X)
% \Picardstack_{X/B} denotes the Picard stack of X over B
% \Picardfunctor_{X/B} denotes the Picard functor of X over B
\def\Pic{\mathop{\mathrm{Pic}}\nolimits}
\def\Picardstack{\mathcal{P}\!\mathit{ic}}
\def\Picardfunctor{\mathrm{Pic}}
\def\Deformationcategory{\mathcal{D}\!\mathit{ef}}
\usepackage{cleveref}
\crefname{originalreference}{original reference}{original references}
\Crefname{originalreference}{Original reference}{Original references}
\makeatletter
\newlabel{sites-modules-definition-site-local}{{03DL}{source}{}{sourcepointer-1}{}}
\newlabel{sites-modules-definition-site-local@cref}{{[originalreference][0][]03DL}{[1][0][]source}{}{}{}}
\makeatother
\crefname{originalequation}{original equation}{original equations}
\Crefname{originalequation}{Original equation}{Original equations}
\makeatletter
\newlabel{equation-alternative-small-ringed}{{0H03}{source}{}{sourcepointer-2}{}}
\newlabel{equation-alternative-small-ringed@cref}{{[originalequation][0][]0H03}{[1][0][]source}{}{}{}}
\makeatother
\crefname{originalreference}{original reference}{original references}
\Crefname{originalreference}{Original reference}{Original references}
\makeatletter
\newlabel{descent-lemma-standard-covering-Cech}{{03FI}{source}{}{sourcepointer-3}{}}
\newlabel{descent-lemma-standard-covering-Cech@cref}{{[originalreference][0][]03FI}{[1][0][]source}{}{}{}}
\makeatother
\crefname{originalreference}{original reference}{original references}
\Crefname{originalreference}{Original reference}{Original references}
\makeatletter
\newlabel{descent-lemma-ff-exact}{{023M}{source}{}{sourcepointer-4}{}}
\newlabel{descent-lemma-ff-exact@cref}{{[originalreference][0][]023M}{[1][0][]source}{}{}{}}
\makeatother
\crefname{originalreference}{original reference}{original references}
\Crefname{originalreference}{Original reference}{Original references}
\makeatletter
\newlabel{sites-modules-lemma-special-locally-free}{{03DM}{source}{}{sourcepointer-5}{}}
\newlabel{sites-modules-lemma-special-locally-free@cref}{{[originalreference][0][]03DM}{[1][0][]source}{}{}{}}
\makeatother
\crefname{originalreference}{original reference}{original references}
\Crefname{originalreference}{Original reference}{Original references}
\makeatletter
\newlabel{schemes-lemma-widetilde-pullback}{{01I9}{source}{}{sourcepointer-6}{}}
\newlabel{schemes-lemma-widetilde-pullback@cref}{{[originalreference][0][]01I9}{[1][0][]source}{}{}{}}
\makeatother
\title{Tensor-compatibility characterizes quasi-coherent modules on the affine \textbackslash{}'etale site}
\author{The Stacks Project Authors}
\date{stacks-project:spaces-descent.tex:local-snapshot-20261012; GNU FDL 1.2 or later}
\begin{document}
\maketitle
\noindent\textit{Editorial note.} All delivered mathematical text is the author's own, copied verbatim from the pinned Stacks Project chapter file spaces-descent.tex; no statement or proof step was rewritten, reordered or removed.\par\medskip
\noindent\textit{Editorial note.} The counted span cites two Descent lemmas, one Modules on Sites lemma and one Schemes lemma, recorded in assembly.external\_references with canonical tags as navigation pointers.\par\medskip
\noindent\textit{Allowed standard prerequisites (editorial).}\begin{itemize}
\item Faithful flatness of standard covering \$\textbackslash{}check\{\textbackslash{}mathrm\{C\}\}\$ech algebra. Stacks Project, Descent, Lemma descent-lemma-standard-covering-Cech. Shows \$\textbackslash{}mathcal\{O\}\_X(U) \textbackslash{}to \textbackslash{}prod \textbackslash{}mathcal\{O\}\_X(U\_i)\$ is faithfully flat.
\item Exactness after faithful flat base change. Stacks Project, Descent, Lemma descent-lemma-ff-exact. Gives exactness of the tensor-amplified sequence.
\item Intrinsicness of quasi-coherent modules on ringed sites. Stacks Project, Modules on Sites, Lemma sites-modules-lemma-special-locally-free. Supports the (2)-(3) transfer along the site equivalence.
\item Tensor of sections with pullback. Stacks Project, Schemes, Lemma schemes-lemma-widetilde-pullback. Concludes (1) from (3).
\end{itemize}\medskip
\ifdefined\hypertarget\hypertarget{sourcepointer-1}{}\fi
\noindent\textit{Source reference.} Original reference 03DL: sources/stacks-chow/tags.txt, source lines 4424--4424. Cross-chapter standard prerequisite cited by the author inside the counted span; it is used as a pointer and its own proof is not part of this item.\par\medskip
\ifdefined\hypertarget\hypertarget{sourcepointer-2}{}\fi
\noindent\textit{Source reference.} Original equation 0H03: sources/stacks-chow/tags.txt, source lines 21198--21198. Cross-chapter standard prerequisite cited by the author inside the counted span; it is used as a pointer and its own proof is not part of this item.\par\medskip
\ifdefined\hypertarget\hypertarget{sourcepointer-3}{}\fi
\noindent\textit{Source reference.} Original reference 03FI: sources/stacks-chow/tags.txt, source lines 4493--4493. Cross-chapter standard prerequisite cited by the author inside the counted span; it is used as a pointer and its own proof is not part of this item.\par\medskip
\ifdefined\hypertarget\hypertarget{sourcepointer-4}{}\fi
\noindent\textit{Source reference.} Original reference 023M: sources/stacks-chow/tags.txt, source lines 2747--2747. Cross-chapter standard prerequisite cited by the author inside the counted span; it is used as a pointer and its own proof is not part of this item.\par\medskip
\ifdefined\hypertarget\hypertarget{sourcepointer-5}{}\fi
\noindent\textit{Source reference.} Original reference 03DM: sources/stacks-chow/tags.txt, source lines 4425--4425. Cross-chapter standard prerequisite cited by the author inside the counted span; it is used as a pointer and its own proof is not part of this item.\par\medskip
\ifdefined\hypertarget\hypertarget{sourcepointer-6}{}\fi
\noindent\textit{Source reference.} Original reference 01I9: sources/stacks-chow/tags.txt, source lines 1970--1970. Cross-chapter standard prerequisite cited by the author inside the counted span; it is used as a pointer and its own proof is not part of this item.\par\medskip
\begin{lemma}
\label{lemma-quasi-coherent-alternative-small}
Let $S$ be a scheme. Let $X$ be an algebraic space over $S$.
Let $\mathcal{F}$ be a presheaf of $\mathcal{O}_X$-modules
on $X_{affine, \etale}$. The following are equivalent
\begin{enumerate}
\item for every morphism $U \to U'$ of $X_{affine, \etale}$ the map
$\mathcal{F}(U') \otimes_{\mathcal{O}_X(U')} \mathcal{O}_X(U)
\to \mathcal{F}(U)$ is an isomorphism,
\item $\mathcal{F}$ is a quasi-coherent module on the ringed site
$(X_{affine, \etale}, \mathcal{O}_X)$ in the sense of
Modules on Sites, Definition \ref{sites-modules-definition-site-local},
\item $\mathcal{F}$ corresponds to a quasi-coherent module on $X$
via the equivalence (\ref{equation-alternative-small-ringed}),
\end{enumerate}
\end{lemma}

\begin{proof}
Assume (1) holds. To show that $\mathcal{F}$ is a sheaf, let
$\mathcal{U} = \{U_i \to U\}_{i = 1, \ldots, n}$ be a covering
of $X_{affine, \etale}$. The sheaf condition for $\mathcal{F}$
and $\mathcal{U}$, by our assumption on $\mathcal{F}$,
reduces to showing that
$$
0 \to \mathcal{F}(U) \to
\prod \mathcal{F}(U) \otimes_{\mathcal{O}_X(U)} \mathcal{O}_X(U_i) \to
\prod \mathcal{F}(U) \otimes_{\mathcal{O}_X(U)} \mathcal{O}_X(U_i \times_U U_j)
$$
is exact. This is true because $\mathcal{O}_X(U) \to \prod \mathcal{O}_X(U_i)$
is faithfully flat
(by Descent, Lemma \ref{descent-lemma-standard-covering-Cech} and
the fact that coverings in $X_{affine, \etale}$ are standard \'etale
coverings) and we may apply Descent, Lemma \ref{descent-lemma-ff-exact}.
Next, we show that $\mathcal{F}$ is quasi-coherent on $X_{affine, \etale}$.
Namely, for $U$ in $X_{affine, \etale}$, set $R = \mathcal{O}_X(U)$
and choose a presentation
$$
\bigoplus\nolimits_{k \in K} R
\longrightarrow
\bigoplus\nolimits_{l \in L} R
\longrightarrow
\mathcal{F}(U)
\longrightarrow 0
$$
by free $R$-modules. By property (1) and the right exactness of tensor product
we see that for every morphism $U' \to U$ in $X_{affine, \etale}$
we obtain a presentation
$$
\bigoplus\nolimits_{k \in K} \mathcal{O}_X(U')
\longrightarrow
\bigoplus\nolimits_{l \in L} \mathcal{O}_X(U')
\longrightarrow
\mathcal{F}(U')
\longrightarrow 0
$$
In other words, we see that the restriction of $\mathcal{F}$
to the localized category $X_{affine, etale}/U$ has a presentation
$$
\bigoplus\nolimits_{k \in K} \mathcal{O}_X|_{X_{affine, \etale}/U}
\longrightarrow
\bigoplus\nolimits_{l \in L} \mathcal{O}_X|_{X_{affine, \etale}/U}
\longrightarrow
\mathcal{F}|_{X_{affine, \etale}/U}
\longrightarrow 0
$$
as required to show that $\mathcal{F}$ is quasi-coherent.
With apologies for the horrible notation, this finishes the proof
that (1) implies (2).

\medskip\noindent
Since the notion of a quasi-coherent module is intrinsic
(Modules on Sites, Lemma \ref{sites-modules-lemma-special-locally-free})
we see that the equivalence (\ref{equation-alternative-small-ringed})
induces an equivalence between categories of quasi-coherent modules.
Thus we have the equivalence of (2) and (3).

\medskip\noindent
Let us assume (3) and prove (1). Namely, let
$\mathcal{G}$ be a quasi-coherent module on $X$ corresponding to $\mathcal{F}$.
Let $h : U \to U' \to X$ be a morphism of $X_{affine, \etale}$.
Denote $f : U \to X$ and $f' : U' \to X$ the structure morphisms,
so that $f = f' \circ h$. We have
$\mathcal{F}(U') = \Gamma(U', (f')^*\mathcal{G})$ and
$\mathcal{F}(U) = \Gamma(U, f^*\mathcal{G}) = \Gamma(U, h^*(f')^*\mathcal{G})$.
Hence (1) holds by Schemes, Lemma \ref{schemes-lemma-widetilde-pullback}.
\end{proof}

\end{document}
